\section{Discussion and Conclusion}\label{sec:discussion}

This section first discusses some potential future improvements to the Isabelle-LLVM based Pastèque implementation in \refsec{sec:future-development}.
\refsec{sec:pacheck-comparison} shows some differences between Pacheck and Pastèque and discusses whether approaches used in Pacheck could also be realized in Pastèque.
\refsec{sec:ai-usage-discussion} and \refsec{sec:takeaways} then close the thesis by a discussion on the use of generative AI tools and my subjective experience of working with Isabelle and Isabelle-LLVM.

\subsection{Further Improvements to Pastèque}\label{sec:future-development}
This subsection discusses some unordered potential improvements to the Isabelle-LLVM migration which might further improve its performance and/or code quality.

Using the currently developed fractional separation logic for Isabelle-LLVM \cite{isabelle-llvm-fractional} in Pastèque might obsolete my custom workarounds of (interleaving) fold-based read only operations on lists (cf. \refsec{sec:non-destructive-list-access-using-fold}, \refsec{sec:interleaving-fold}) and the fused addition (cf. \refsec{sec:avoiding-polynomial-deep-copies}) while also making it possible to further avoid deep-copy operations.

Another optimization, which could already be carried out without fractional separation logic, might be to represent terms in monomials by arrays instead of lists, as is currently the case.
This optimization might speed up comparison between terms, allowing to shortcut comparison when array lengths are unequal and would certainly reduce the total number of allocations.

Furthermore, the current design of partial maps (cf. \refsec{sec:low-level-refinement-of-lists-of-optionals}) wraps the contained elements (in case of Pastèque polynomials) by a pointer to obtain the null pointer as a representation of empty slots.
The contained polynomials are already represented by a pointer (which is the refinement target of the outer list), so effectively this introduces a double pointer dereference on lookup.
It might be possible to instead define the null pointer as the explicit refinement target of the $\iexp{None}$ case of the high-level $\iexp{sllist\_polynomial} \app \iexp{option}$ type to save the additional pointer dereference.

Another possible change to Pastèque would be to move parsing to the Isabelle side such that the C side only handles IO and hands the files directly as strings to the verified LLVM code.
While it is unlikely that this change would improve performance, it would reduce the amount of trusted parsing code and allow for a verified parsing process.
However, one would have to formally define a semantic for the given input files to obtain a meaningful correctness proof which is a non-trivial task.

\subsection{Comparison to Pacheck}\label{sec:pacheck-comparison}
The evaluation in \refsec{sec:evaluation} shows that Pacheck is faster than both variants of Pastèque.
This subsection notes some differences of Pacheck, compared to Pastèque, and discusses whether a similar approach could possibly be used for Pastèque.

Variables (i.e. strings) in Pacheck are unique; if a variable occurs in multiple terms, the terms will reference the same variable object.
This is also implemented in Pastèque through the shared variable optimization, given by $\iexp{NV}$ and $\iexp{VN}$ (cf. \refsec{sec:methodology_high-level-representation-of-polynomials}).
Like Pastèque, Pacheck implements terms as linked lists of variables.
Other that in Pastèque however, terms are unique in Pacheck; if two terms have a common suffix, they will reference the same chain of nodes.
In Pasteque, two terms would be two separate lists, as indicated in Figure~\ref{fig:pacheck-shared-suffix}.
\begin{figure}[htbp]
	\centering
	\begin{tikzpicture}[
			var/.style={draw, rectangle, minimum width=7mm, minimum height=7mm, font=\ttfamily, text height=1.5ex, text depth=.25ex},
			heading/.style={text height=2ex, text depth=.25ex},
			ptr/.style={->, >=stealth, thick},
			node distance=6mm and 6mm,
		]
		% Pastèque: one separate list per term
		\node[var] (pa) {a};
		\node[var, right=of pa] (pb1) {b};
		\node[var, right=of pb1] (pc1) {c};
		\node[var, right=of pc1] (pd1) {d};
		\node[var, below=of pa] (px) {x};
		\node[var, right=of px] (pb2) {b};
		\node[var, right=of pb2] (pc2) {c};
		\node[var, right=of pc2] (pd2) {d};
		\draw[ptr] (pa) -- (pb1);
		\draw[ptr] (pb1) -- (pc1);
		\draw[ptr] (pc1) -- (pd1);
		\draw[ptr] (px) -- (pb2);
		\draw[ptr] (pb2) -- (pc2);
		\draw[ptr] (pc2) -- (pd2);
		\node[heading, above=4mm] at ($(pa.north)!0.5!(pd1.north)$) {Pastèque};

		% Pacheck: terms share their common suffix
		\node[var, right=18mm of pd1] (ka) {a};
		\node[var, right=of ka] (kb) {b};
		\node[var, right=of kb] (kc) {c};
		\node[var, right=of kc] (kd) {d};
		\node[var, below=of ka] (kx) {x};
		\draw[ptr] (ka) -- (kb);
		\draw[ptr] (kb) -- (kc);
		\draw[ptr] (kc) -- (kd);
		\draw[ptr] (kx) -- (kb);
		\node[heading, above=4mm] at ($(ka.north)!0.5!(kd.north)$) {Pacheck};
	\end{tikzpicture}
	\caption{Representation of the terms $abcd$ and $xbcd$ as linked lists of variables.
		Pastèque stores both terms as separate lists, whereas in Pacheck the term $xbcd$ references the shared suffix $bcd$ of the term $abcd$.}
	\label{fig:pacheck-shared-suffix}
\end{figure}
The same layout would not be possible in Pastèque as refinement assertions for the two linked lists could not reference the same suffix.
Consequently, copying terms or monomials in Pacheck amounts to increasing their reference counter, in Pastèque, copying actually creates a copy of the respective object in memory.
Checking equality of terms in Pacheck amounts to a pointer comparison, while in Pastèque, all pairs of variables are compared until an unequal pair is found.

Pacheck interleaves parsing and checking, similar to the \inlinecode{uloop} option in Pastèque (which only exists for the Standard ML backend).
If parsing was moved from C to Isabelle, it would be possible to also interleave parsing and checking in Pastèque while still fully verifying correctness.

Lastly, Pacheck uses GMP \cite{gmp} for multiple precision arithmetic which is likely faster than the currently used custom multiple precision integer library by Spinei \cite{spinei_isabelle_big_integer} in the LLVM backend of Pastèque.
As of now, it is unclear however how much of an impact this has on the performance difference between Pastèque and Pacheck.

\subsection{AI-Assisted Proof Development}\label{sec:ai-usage-discussion}
When starting to work on this thesis, I initially set a goal of not using any generative AI tools for Isabelle proofs or implementation.
For the initial work of extending the multiple precision integer arithmetic library, discussed in \refsec{sec:multiple-precision-integers}, I had followed that goal and did not use generative AI tools \footnote{For the final implementation of conversion between integer and string, I did use AI to write parts of the proofs.}.
Once I moved onto the Isabelle-LLVM migration and realized that I would have to implement fundamental low-level data structures, discussed in \refsec{sec:low-level-datatypes-in-isabelle-llvm}, I gave up on the goal of not using AI tools after failing to proof correct even simple fundamental operations by myself.
I then started to use the tool Claude Code by Anthropic, using different versions of the language models "Opus", "Sonnet" and "Fable", by copying proof goals into Claude Code and task it to analyze the goal and draft initial proofs, which I would then copy back into the actual Isabelle theories.
Often, the generated proofs were not accepted by the Isabelle kernel directly and I could use Isabelle's built in proof automation, in particular sledgehammer, to discharge failing steps in the proof.
Some parts of the final theories, in particular the theory $\iexp{LLVM\_DS\_Char\_Literals}$ which contains refinement of printable ASCII characters, were directly generated by Claude Code.

Compared to using a web-based chat interface, using a local installation of an "agentic" tool like Claude Code has the advantage that the tool will use command line tools like \inlinecode{grep} to add relevant definitions from the Isabelle-LLVM library to its context window which, in my experience, improves the quality of its outputs significantly and makes it a viable tool for this use case.

\subsection{Working with Isabelle-LLVM}\label{sec:takeaways}
While some particular low-level optimizations like bit-stuffing are not possible in Isabelle-LLVM, it still moves the formally verified approach significantly closer to high-performance languages by enabling simultaneous development and verification of imperative programs with manual memory management while also allowing for elegant correctness proofs through its step-wise refinement approach.

I am convinced that my work on the Pastèque migration to Isabelle-LLVM is not representative of Isabelle-LLVM's full potential regarding efficiency and that it would be possible to improve the LLVM migration to supersede the performance of the Standard ML implementation of Pastèque in 512-bit instances aswell.

My personal biggest struggles of working with Isabelle-LLVM mostly stem from lack of introductory documentation, i.e., a "beginner tutorial" specific to Isabelle-LLVM, extending the quickstart guide of the imperative refinement framework \cite{Refine_Imperative_HOL-AFP} with information on how to work with $\iexp{llM}$-level functions directly.
Another hurdle is intrinsic to its approach of being implemented as proof methods for Isabelle; when refining functions with \inlinecode{sepref} or proving $\iexp{llM}$-level theorems with \inlinecode{vcg}, there are no error messages when something does not go as expected.
Instead, the tools provide feedback through unresolved goal states which are often lengthy and contain symbols and definitions from the underlying separation logic, which are usually hidden from the user when writing theorems.
To productively work with Isabelle-LLVM, it takes significant amounts of time to obtain enough intuition to interpret the returned goal states.
In particular the \inlinecode{vcg} tool might run into unprovable goal states as a result of over eager simplifications by the user which I often could not detect.
I had to resort to using Claude Code to overcome such situations, which I originally wanted to avoid.

With that being said, I found Isabelle-LLVM's interactive approach to formal verification of complex software, and working with Isabelle in general to be immensely insightful, being a hands-on combination of theory and practice that I have not seen anywhere else.
